\begin{frame}{자주 묻는 질문 1: 왜 세 가지로 나뉘나}
  \kq{Q·K·V는 왜 같은 입력에서 서로 다른 세 가지로 나뉘나요?
  입력 자체를 비교하면 안 되나요?}
  \begin{itemize}
    \item $Q=K=V=X$ (변환 없이 원본 입력을 그대로 씀)라면
      ``이 단어가 자기 자신과 얼마나 비슷한가''만 재게 되어
      \textbf{표현력이 크게 제한}된다.
    \item 서로 다른 학습 가능한 변환 $W_Q, W_K, W_V$ 를 두면,
      ``무엇을 찾을지''와 ``무엇을 제공할지''를 모델이
      \textbf{데이터로부터 독립적으로 학습}할 수 있다.
    \item 예: 대명사는 ``나를 가리키는 명사를 찾는다'' 방향으로
      $W_Q$ 가, 명사는 ``대명사가 찾는 대상이 될 수 있다'' 방향으로
      $W_K$ 가 학습된다.
  \end{itemize}
  \vskip0.5em
  {\footnotesize 보충: ``내적(유사도)으로 관련을 재고, 학습된 변환을
  거친 뒤 비교한다''는 패턴은 Week 5.3 커널 함수 $K(x,x')$ 과
  정신적으로 닮아 있다.}
\end{frame}
